% CERT rel=rec a=1 b=2 f=1 T1=1 cerrada=log2(n)+1
\ej{4}{Búsqueda binaria: sustitución e inducción.}

Consideramos $n=2^k$, con $k\in\N\cup\{0\}$, para que las divisiones entre $2$ alcancen el caso base $T(1)=1$.

Conjeturamos $T(n)\in\Th{\log n}$. Para demostrarlo por sustitución, probaremos por inducción la fórmula más precisa $T(n)=\log_2 n+1$.

\begin{enumerate}
\item \textbf{Caso base ($k=0$).} Como $n=1$, se cumple $T(1)=1=\log_2 1+1$.

\item \textbf{Hipótesis inductiva.} Supongamos que $T(2^k)=k+1$ para algún entero $k\ge0$.

\item \textbf{Paso inductivo.} Aplicando la recurrencia y sustituyendo la hipótesis inductiva:
\[
T(2^{k+1})=T(2^{k+1}/2)+1
=T(2^k)+1=(k+1)+1
=\log_2(2^{k+1})+1.
\]
Por inducción, $T(n)=\log_2 n+1$ para toda potencia de $2$ con $n\ge1$.
\end{enumerate}

Para concluir $T(n)\in\Th{\log n}$, debemos encontrar $c_1,c_2>0$ y $n_0\in\N$ tales que $0\le c_1\log_2 n\le T(n)\le c_2\log_2 n$ para todo $n\ge n_0$ del dominio.

Para $n\ge2$, tenemos $\log_2 n\ge1$; por ello,
\[
0\le\log_2 n\le\log_2 n+1=T(n)
\le\log_2 n+\log_2 n=2\log_2 n.
\]
Podemos tomar $c_1=1$, $c_2=2$ y $n_0=2$. Por tanto, $\boxed{T(n)\in\Th{\log n}}$.
\fin
